% !TEX encoding = UTF-8
% !TEX TS-program = pdflatex
% !TEX root = ../tesi.tex

%**************************************************************
\chapter{Conclusions}
\label{cap:conclusions}
%**************************************************************
{\color{blue} 

This thesis investigated the application of Reinforcement Learning (RL) to the problem of Dynamic Directional Modulation (DDM), with the objective of enabling the real-time generation of secure antenna array configurations that simultaneously satisfy communication and security constraints. Unlike conventional Directional Modulation approaches, which typically rely on deterministic optimization procedures to compute a single beamforming solution, Dynamic Directional Modulation requires the continuous generation of multiple valid configurations in order to prevent an eavesdropper from exploiting the static characteristics of the transmitted signal. This requirement naturally motivates a stochastic learning framework capable of producing a distribution of feasible solutions rather than a single deterministic optimum.

The first contribution of this work is the formulation of Dynamic Directional Modulation as a contextual bandit problem. Within this framework, the desired transmission direction acts as the context, while the antenna  phases constitute the action space. By learning a stochastic policy, the proposed approach is able to generate diverse antenna configurations that satisfy the desired radiation pattern constraints while maintaining the dynamic behaviour required by DDM. The results demonstrate that RL can effectively learn this mapping and perform inference in real time, replacing computationally expensive optimization procedures with a single forward pass through a neural network.

A second contribution is the geometric analysis of the DDM optimization problem. Through the study of the array factor and the phase-flatness constraint, it was shown that the difficulty of the problem is strongly dependent on the target direction. In particular, directions close to broadside were found to correspond to significantly smaller feasible regions in the solution space, making them intrinsically more challenging than directions near endfire. This observation provides an important physical interpretation of the learning dynamics and explains several convergence difficulties encountered during training.

Building upon this analysis, a third contribution of the thesis is the identification of a sampling imbalance introduced by conventional uniform-angle training strategies. Because the array response depends on the spatial-frequency variable $u=\sin(\varphi)$, uniform sampling in the angular domain does not correspond to a uniform coverage of the underlying physical problem. To address this issue, a sampling strategy based on uniform spacing in spatial-frequency space was proposed and integrated into the training process. This modification provides a more balanced representation of the problem domain and improves exploration of the most challenging regions of the angular range.

The fourth contribution is the development of a curriculum-learning framework specifically tailored to Dynamic Directional Modulation. The proposed curriculum progressively increases task complexity by initially focusing on the most difficult broadside directions and gradually expanding the angular range and tightening the design constraints. This strategy was found to significantly improve training stability and convergence, enabling the policy to discover feasible solutions in regions where direct training frequently struggled.

The experimental results illustrqte that the  learned policy is capable of producing distinct antenna excitation patterns for the same transmission direction while maintaining the desired radiation characteristics, thereby realizing the fundamental objective of Dynamic Directional Modulation. The approach performs particularly well for lower-order modulation schemes, while the results obtained for higher-order constellations reveal the increasing complexity of the underlying optimization landscape.

To sum-up, this thesis demonstrates that Reinforcement Learning provides a viable and flexible framework for Dynamic Directional Modulation. Beyond the successful application of RL itself, the work contributes new insights into the geometry of the DDM problem, highlights the importance of spatial-frequency-based sampling, and introduces curriculum-learning strategies that make the problem tractable in practice. These results establish a foundation for future research at the intersection of machine learning, antenna array processing, and physical-layer security.}

The proposed framework demonstrates promising results for lower modulation orders, successfully steering the array and maintaining acceptable symbol error rates under the simulated conditions. However, scalability remains a limiting factor: performance degrades notably at higher modulation orders such as 16-PSK, where the tighter noise margin amplifies the sensitivity to suboptimal beamforming configurations.
The central property of the framework is the stochasticity introduced at inference time: rather than a single deterministic output or a fixed beamforming configuration, the model generates a distribution of solutions, realising dynamic directional modulation.
The two training variations explored in this work reinforce this picture: the physics informed parameterisation confirms that embedding analytical structure in the action decoder is viable but reduces solution diversity and limits the policy's use of element-silencing strategies while the $N=8$ experiment shows that scaling the array exacerbates the model's limitations. Even the preferred gradient-based approach would require tuning, such as a larger exploration budget, a revised curriculum, or the use of known valid solutions as initialization points, to maintain the same quality of results achieved at $N=5$ when scaling with higher element numbers. 
The results are promising as a proof of concept, yet the gap between the simulated environment and real-world deployment remains substantial. Bridging this gap constitutes the primary direction for future development, as discussed in \chaptername~\ref{cap:future_works}.

%**************************************************************
\section{Future work}
\label{cap:future_works}
For future development on this thesis's groundwork, several directions to increase physical fidelity present themselves as natural extensions of this work.

\begin{itemize}
    \item \textbf{Hardware impairments and practical array constraints:} The current model assumes continuous-valued phase shifters and ideal, independent radiating elements; future works should incorporate finite phase-resolution constraints into the environment, alongside additional impairments such as nonlinear amplifiers, in-phase/quadrature imbalance, and oscillator noise. These effects introduce systematic deviations between the idealized steering vector and the physically realized beam pattern.
    \item \textbf{Multi-agent and distributed beamforming:} For very large arrays, such as those utilized in massive MIMO, centralized control of all antenna weights becomes computational prohibitive. A multi-agent reinforcement learning (MARL) formulation, in which individual agents govern sub-array or antenna panels and coordinate through limited inter-agent communication, offers a scalable alternative. 
    \item \textbf{Uniform Planar Array and full three-dimensional steering:} The current environment is restricted to a Uniform Linear Array (ULA) operating in the azimuth plane at a fixed elevation of $theta = 0^\circ$. Extending the geometry to a Uniform Planar Array (UPA) introduces a two-dimensional action space over both azimuth $\varphi$ and elevation $\theta$, increasing the dimensionality of the optimization problem and requiring the agent to learn a richer mapping from angular coordinates to beamforming weights.
    \item \textbf{Realistic channel models and multi-path propagation:} The reward signal currently evaluates beam quality against a free-space steering vector, which assumes a single line-of-sight propagation path. Practical wireless channels exhibit multi-path fading and Doppler shift. Replacing the idealized channel model with a standardized stochastic model would expose the agent to the full distribution of channel realizations it would encounter in real-world applications, promoting robustness over a wider range of propagation conditions. 
\end{itemize}

%**************************************************************
